\documentclass[11pt]{article}
\usepackage[margin=1in]{geometry}
\usepackage[T1]{fontenc}
\usepackage[utf8]{inputenc}
\usepackage{hyperref}
\usepackage{graphicx}
\usepackage{amsmath}
\usepackage{amssymb}
\usepackage{booktabs}
\usepackage{listings}
\usepackage{xcolor}

\title{Correctness and Token-Per-Second Optimization for LLaMA-Style Inference in Raw Rust}
\author{Student Notes (generated for learning)}
\date{February 21, 2026}

\lstdefinelanguage{Rust}{
  keywords={fn,let,mut,struct,enum,impl,trait,for,while,loop,if,else,match,return,Self,crate,super,mod,pub,use,where,async,await},
  keywordstyle=\color{blue}\bfseries,
  comment=[l]{//},
  commentstyle=\color{gray},
  stringstyle=\color{orange},
  morestring=[b]",
}

\lstset{
  language=Rust,
  basicstyle=\ttfamily\small,
  breaklines=true,
  columns=fullflexible,
}

\begin{document}
\maketitle

\begin{abstract}
This document is a practical, student-friendly guide to improving correctness and throughput (tokens per second, or tok/s) for LLaMA-style inference engines implemented in ``raw'' Rust. It aims to explain the engineering processes involved in building and validating a transformer inference stack from scratch, while also showing how performance improvements on consumer hardware are discovered, measured, and kept correct. The focus is on: (1) correctness across numerical precision, attention implementations, and sampling, and (2) performance across memory layout, kernel fusion, and threading.
\end{abstract}

\tableofcontents
\newpage

\section{Context and Goals}
You have a Rust-based LLaMA implementation. You want it to be:
\begin{itemize}
  \item \textbf{Correct}: matches reference outputs within expected tolerance.
  \item \textbf{Fast}: competitive tok/s on consumer CPUs/GPUs.
  \item \textbf{Understandable}: you can explain why it works.
\end{itemize}

These goals often conflict, so the work is a balancing act between math, systems programming, and scientific testing.

\section{What ``Correctness'' Means in LLM Inference}
Correctness does \emph{not} mean bit-identical output across all backends. Floating point math, parallelism, and fused kernels introduce small differences. A good correctness strategy uses multiple layers:

\subsection{Layer 0: Sanity Checks}
\begin{itemize}
  \item Deterministic decoding with a fixed seed.
  \item Known prompts with known top-1 token from a reference implementation.
  \item Small models (e.g., tiny LLaMA) where you can run exhaustive comparisons.
\end{itemize}

\subsection{Layer 1: Unit-Level Equivalence}
Validate each subcomponent in isolation:
\begin{itemize}
  \item Linear layers vs. reference: $Y = XW^T + b$.
  \item RMSNorm vs. reference.
  \item RoPE (rotary positional embedding) correctness.
  \item Attention math: masking, scaling, and softmax.
\end{itemize}

\subsection{Layer 2: End-to-End Equivalence}
Compare against a trusted backend (e.g., PyTorch) with tolerances:
\begin{itemize}
  \item Same weights, same prompt, same decoding, same seed.
  \item Validate logits at each step.
  \item Use $L_\infty$ or $L_2$ error thresholds per layer or per token.
\end{itemize}

\subsection{Layer 3: Behavioral Correctness}
When exact numeric equivalence fails, use behavioral tests:
\begin{itemize}
  \item Perplexity over a small dataset.
  \item Top-k overlap (how often do top-k tokens match).
  \item Sampling distribution similarity.
\end{itemize}

\section{Floating Point Realities}
Rust, like C++, does not change floating point behavior by itself. Differences come from:
\begin{itemize}
  \item \textbf{Precision}: f32 vs f16 vs bf16.
  \item \textbf{Order of operations}: parallel reduction order changes rounding.
  \item \textbf{Fused kernels}: combining operations alters rounding.
  \item \textbf{Vectorization}: SIMD instructions are not bitwise deterministic.
\end{itemize}

Therefore, set up tolerances and test plans accordingly.

\section{A Minimal LLaMA Inference Pipeline}
A simplified pipeline per token:
\begin{enumerate}
  \item Embed input token.
  \item For each transformer layer:
    \begin{itemize}
      \item RMSNorm on input.
      \item QKV projection.
      \item RoPE on Q and K.
      \item Attention: compute scores, apply mask, softmax, weighted sum.
      \item Output projection.
      \item MLP: gated activation + output projection.
      \item Residual adds.
    \end{itemize}
  \item Final RMSNorm and output projection to logits.
  \item Sampling: greedy, top-k, top-p, temperature.
\end{enumerate}

Each bullet point is a place where correctness can break or performance can be optimized.

\section{Correctness Pitfalls (Common Bugs)}
\begin{itemize}
  \item \textbf{Wrong shape assumptions}: attention head dimension mismatch.
  \item \textbf{RoPE offsets}: incorrect position index when caching.
  \item \textbf{KV cache layout bugs}: misaligned cache strides.
  \item \textbf{Masking bugs}: causal mask not applied correctly for new tokens.
  \item \textbf{Sampling mistakes}: softmax over wrong dimension or wrong temperature.
  \item \textbf{Activation mismatch}: using GELU instead of SwiGLU.
\end{itemize}

\section{KV Cache: The Center of Both Correctness and Speed}
The KV cache stores keys and values for previous tokens. For each new token:
\begin{itemize}
  \item Compute K and V for that token.
  \item Append into the cache.
  \item Use cached K/V for attention scores.
\end{itemize}

\subsection{Correctness Concerns}
\begin{itemize}
  \item Position index must match cache slot.
  \item Cache layout must be consistent across layers and heads.
  \item Any quantized cache needs careful dequantization.
\end{itemize}

\subsection{Performance Concerns}
\begin{itemize}
  \item Cache layout should be contiguous and cache-friendly.
  \item Avoid per-token allocations.
  \item Align buffers for SIMD.
\end{itemize}

\section{Measuring tok/s}
Tok/s is the main performance metric for inference. Basic formula:
\[\text{tok/s} = \frac{\text{tokens decoded}}{\text{seconds elapsed}}\]

Best practice:
\begin{itemize}
  \item Separate \textbf{prefill} (processing the prompt) from \textbf{decode} (one token at a time).
  \item Measure both:\; prefill tok/s and decode tok/s.
  \item Use multiple runs to average.
\end{itemize}

\section{Optimization Process (High-Level)}
A practical loop:
\begin{enumerate}
  \item Establish a reference test (correctness + baseline tok/s).
  \item Change one thing at a time.
  \item Validate correctness at multiple levels.
  \item Measure tok/s and memory.
  \item Keep the change only if correctness holds and performance improves.
\end{enumerate}

\section{Performance Bottlenecks in Raw Rust}
\subsection{Matrix Multiply (GEMM)}
\begin{itemize}
  \item This dominates runtime.
  \item Use high-quality kernels or SIMD intrinsics.
  \item Pay attention to layout: row-major vs column-major.
\end{itemize}

\subsection{Attention Softmax}
\begin{itemize}
  \item Softmax must be numerically stable.
  \item Compute max, subtract, exponentiate, sum, divide.
  \item For speed, fuse steps if possible.
\end{itemize}

\subsection{Memory Bandwidth}
\begin{itemize}
  \item Large models are memory-bound.
  \item Reuse buffers; avoid copies.
  \item Keep hot data in cache: KV cache layout matters.
\end{itemize}

\section{Typical Rust-Level Optimizations}
\begin{itemize}
  \item \textbf{Avoid bounds checks in hot loops}: use iterators carefully, or unsafe blocks with care.
  \item \textbf{SIMD} with `std::simd` or platform intrinsics.
  \item \textbf{Memory alignment}: use \texttt{\#[repr(align(64))]} for hot buffers.
  \item \textbf{Threading}: rayon or custom thread pools.
  \item \textbf{Minimize allocations}: preallocate scratch buffers.
  \item \textbf{Fused ops}: combine RMSNorm + linear if possible.
\end{itemize}

\section{Quantization and Its Impact}
Quantization compresses weights to 8-bit or 4-bit for faster inference.
\begin{itemize}
  \item Increases tok/s by reducing memory bandwidth.
  \item Risks correctness if scale/zero-point handling is wrong.
  \item Use per-channel scales when possible.
\end{itemize}

\section{Sampling Correctness}
Sampling bugs are easy to miss.
\begin{itemize}
  \item Ensure logits are in float and stable before softmax.
  \item Top-k or top-p should renormalize the distribution.
  \item Temperature scaling must be applied before softmax.
\end{itemize}

\section{Process: What ``Vibe-Coded in 9 Days'' Actually Means}
If you built an implementation quickly, you likely:
\begin{itemize}
  \item Glued together minimal components.
  \item Tested only end-to-end behavior.
  \item Used small prompts and simple sampling.
\end{itemize}

The next step is \emph{making it rigorous}:
\begin{itemize}
  \item Add unit tests per layer.
  \item Add reference comparisons.
  \item Add benchmarks with fixed settings.
  \item Document every assumption.
\end{itemize}

\section{A Concrete Debugging Checklist}
\begin{enumerate}
  \item Compare single-layer outputs to reference.
  \item Validate attention scores before softmax.
  \item Check RoPE against a known implementation.
  \item Ensure cache positions align with token index.
  \item Validate logits for the first 5 tokens.
  \item Measure tok/s before and after changes.
\end{enumerate}

\section{Example: Stable Softmax in Rust}
\begin{lstlisting}
fn softmax(x: &mut [f32]) {
    let mut max = f32::NEG_INFINITY;
    for &v in x.iter() { if v > max { max = v; } }

    let mut sum = 0.0f32;
    for v in x.iter_mut() {
        *v = (*v - max).exp();
        sum += *v;
    }

    let inv = 1.0f32 / sum;
    for v in x.iter_mut() { *v *= inv; }
}
\end{lstlisting}

\section{Testing Strategy (Minimal but Useful)}
\begin{itemize}
  \item \textbf{Golden tests}: known prompt, known logits for a small model.
  \item \textbf{Diff tests}: compare against PyTorch (tolerances).
  \item \textbf{Performance tests}: fixed prompt length, fixed decode length.
  \item \textbf{Regression tests}: run on every change.
\end{itemize}

\section{Hardware Considerations}
Consumer hardware matters:
\begin{itemize}
  \item CPU: number of cores, AVX2/AVX-512 support.
  \item GPU: memory size and bandwidth.
  \item RAM: large models can page if memory is low.
\end{itemize}

You must measure on the actual hardware you care about.

\section{Putting It Together: A Practical Workflow}
\begin{enumerate}
  \item Start with a known-correct model and small prompt.
  \item Implement one component at a time.
  \item Validate each component in isolation.
  \item Integrate into end-to-end.
  \item Optimize the slowest parts only.
  \item Re-test correctness after every optimization.
\end{enumerate}

\section{Appendix: Suggested Reading Topics}
\begin{itemize}
  \item Transformer attention math.
  \item Numerical stability in softmax.
  \item SIMD in Rust.
  \item Quantization fundamentals.
  \item Profiling tools for Rust.
\end{itemize}

\end{document}
